\documentclass[10pt,a4paper]{article}

\usepackage[top=2.0cm, bottom=2.0cm, left=2.0cm, right=2.0cm]{geometry}
\usepackage{amsmath,amssymb,amsfonts}
\usepackage{mathptmx} % Clean Times Roman academic typography
\usepackage{microtype}
\usepackage{booktabs}
\usepackage{tabularx}
\usepackage{array}
\usepackage{titlesec}
\usepackage{fancyhdr}
\usepackage{setspace}
\usepackage[colorlinks=true,linkcolor=black,citecolor=black,urlcolor=black]{hyperref}

% Scholarly single-spaced paragraph formatting
\setstretch{1.10}

\pagestyle{fancy}
\fancyhf{}
\fancyhead[L]{\small\textit{M.Tech Thesis Progress Report --- Work Done Synthesis}}
\fancyhead[R]{\small\textit{P. K. Pandey (SC25M147)}}
\fancyfoot[C]{\thepage}
\renewcommand{\headrulewidth}{0.4pt}
\renewcommand{\footrulewidth}{0pt}

\titleformat{\section}{\large\bfseries}{\thesection.}{0.7em}{}[\titlerule]
\titleformat{\subsection}{\normalsize\bfseries}{\thesubsection}{0.5em}{}
\titleformat{\subsubsection}{\small\bfseries\itshape}{\thesubsubsection}{0.4em}{}
\titlespacing*{\section}{0pt}{1.2em}{0.5em}
\titlespacing*{\subsection}{0pt}{1.0em}{0.4em}
\titlespacing*{\subsubsection}{0pt}{0.8em}{0.3em}

\begin{document}

\begin{center}
    {\LARGE\bfseries Multiphysics Modeling, Solidification Dynamics, and Ultrasonic Coupling in Wire Arc Additive Manufacturing (WAAM)\par}
    \vspace{0.6em}
    {\large\textbf{Comprehensive Technical Synthesis of Research Work Completed (August--October 2026)}\par}
    \vspace{0.6em}
    {\normalsize\textbf{Author:} Pranay K. Pandey (\texttt{SC25M147})\par}
    {\small M.Tech in Manufacturing Technology, Department of Aerospace Engineering\par}
    {\small Indian Institute of Space Science and Technology (IIST), Thiruvananthapuram\par}
    \vspace{0.3em}
    {\small\textbf{Thesis Supervisors:} Dr. Sooraj V. S. \& Dr. Chakravarthy P.\par}
    \vspace{0.2em}
    {\small\textit{Presentation to the Thesis Review Committee / Comprehensive Examination Panel}\par}
\end{center}

\vspace{0.3em}
\hrule
\vspace{0.8em}

\begin{abstract}
\noindent Over the past two months, a rigorous, multiscale computational and experimental research campaign was executed to model and optimize Wire Arc Additive Manufacturing (WAAM) and Ultrasonic-Assisted WAAM (UVA-WAAM). This technical report synthesizes all research milestones achieved across five interrelated engineering domains: (i) an open-source, conforming-mesh, thermomechanical finite-element framework implemented in FreeFEM++ with radial return \(J_2\) plasticity; (ii) high-fidelity Volume-of-Fluid (VOF) weld-pool hydrodynamics developed in OpenFOAM (\texttt{waamFoam}) incorporating Marangoni thermal-solutal convection and non-Boussinesq drag; (iii) a mesoscale dendritic phase-field suite resolving thin-interface solute trapping, microsegregation, and the Columnar-to-Equiaxed Transition (CET); (iv) ultrasonic acoustic cavitation, radiation pressure, acoustic streaming, and dislocation acoustic softening modeling; and (v) pragmatic thermal modeling and physical machine toolpath commissioning at the IIST laboratory. All underlying mathematical formulations, algorithmic corrections, verification benchmarks, and experimental setups are documented in standard scholarly prose.
\end{abstract}

\section{Introduction and Research Rationale}
Wire Arc Additive Manufacturing (WAAM) presents a highly productive pathway for manufacturing large-scale, buy-to-fly-critical components in aerospace applications. By deploying an electric arc as a concentrated thermal source to fuse wire feedstock layer-upon-layer, deposition rates exceeding \(2.0\text{--}4.0\,\mathrm{kg/h}\) are readily achieved. However, WAAM deposits inherently experience severe thermal gradients, cyclic reheat shocks, coarse epitaxially grown columnar grains, anisotropic yield strengths, and significant residual stress accumulation leading to out-of-plane distortion.

To overcome these metallurgical limitations, ultrasonic-assisted WAAM (UVA-WAAM) introduces high-frequency acoustic vibrations during deposition. While experimental literature demonstrates grain refinement and stress relief, the underlying multi-scale physics remains weakly coupled in existing computational models. Over the initial two-month investigation, our objective has been to build an end-to-end, open, mathematically verifiable multiscale framework. Rather than utilizing commercial black-box codes, every physical domain was audited, verified against analytical solutions or published experimental benchmarks, and unified into an open toolchain.

\section{Macroscale Thermomechanical Finite Element Framework}

\subsection{Governing Continuum Mechanics and Phase Evolution}
A complete continuum mechanics solver was developed in FreeFEM++ version 4.15 and parallelized using the MUMPS direct solver. The transient thermal evolution is governed by the nonlinear heat conduction equation:
\[
\rho(T)\,c_p(T)\,\frac{\partial T}{\partial t} = \nabla \cdot \bigl(k(T)\,\nabla T\bigr) + \dot{q}_{\mathrm{ext}} - \rho(T)\,L_f\,\frac{\partial f_l}{\partial t}
\]
where \(\rho(T)\) is the temperature-dependent density, \(c_p(T)\) is the specific heat capacity, \(k(T)\) is the thermal conductivity, and \(L_f\) represents the latent heat of fusion. The liquid fraction \(f_l(T)\) across the solidus--liquidus interval \([T_{\mathrm{sol}},\,T_{\mathrm{liq}}]\) is resolved via an apparent heat capacity formulation:
\[
c_p^*(T) = c_p(T) + L_f\,\frac{\partial f_l}{\partial T}, \quad \text{with} \quad f_l(T) = \frac{T - T_{\mathrm{sol}}}{T_{\mathrm{liq}} - T_{\mathrm{sol}}}
\]
Thermal boundary conditions enforce both Newton cooling and Stefan-Boltzmann radiation:
\[
-k\,\nabla T \cdot \vec{n} = h_{\mathrm{conv}}\,(T - T_\infty) + \varepsilon\,\sigma_{\mathrm{SB}}\,(T^4 - T_\infty^4)
\]
where \(h_{\mathrm{conv}}\) is the convective heat transfer coefficient, \(\varepsilon\) is surface emissivity, and the Stefan-Boltzmann constant is \(\sigma_{\mathrm{SB}} = 5.67037\times 10^{-8}\,\mathrm{W/(m^2\,K^4)}\).

\subsection{Elastoplasticity and Radial Return Algorithm}
The quasi-static momentum balance neglects inertial effects:
\[
\nabla \cdot \boldsymbol{\sigma} = \vec{0}, \quad \boldsymbol{\varepsilon} = \frac{1}{2}\bigl(\nabla \vec{u} + (\nabla \vec{u})^T\bigr)
\]
Total strain is additively decomposed into elastic, plastic, and thermal components:
\[
\boldsymbol{\varepsilon} = \boldsymbol{\varepsilon}^e + \boldsymbol{\varepsilon}^p + \boldsymbol{\varepsilon}^{\mathrm{th}}, \quad \boldsymbol{\varepsilon}^{\mathrm{th}} = \alpha_{\mathrm{sec}}(T)\,(T - T_{\mathrm{ref}})\,\mathbf{I}
\]
An incremental \(J_2\) von Mises plasticity model with isotropic hardening was formulated. For each load increment, trial elastic stress is evaluated:
\[
\boldsymbol{\sigma}^{\mathrm{trial}} = \mathbf{C}(T) : (\boldsymbol{\varepsilon}_{n+1} - \boldsymbol{\varepsilon}^p_n - \boldsymbol{\varepsilon}^{\mathrm{th}}_{n+1})
\]
Yielding occurs when the trial yield function \(f^{\mathrm{trial}} = \sigma_{\mathrm{eq}}^{\mathrm{trial}} - \sigma_y(T,\,p_n) > 0\). The plastic multiplier increment \(\Delta p\) is computed by solving:
\[
\Delta p = \frac{\sigma_{\mathrm{eq}}^{\mathrm{trial}} - \sigma_y(T,\,p_n)}{3\,G(T) + H'(T)}
\]
where \(G(T)\) is the shear modulus and \(H'\) is the plastic hardening slope. The back-projection of deviatoric stress is exact:
\[
\mathbf{s}_{n+1} = \mathbf{s}^{\mathrm{trial}} - 2\,G(T)\,\Delta p\,\mathbf{n}^{\mathrm{trial}}, \quad \mathbf{n}^{\mathrm{trial}} = \frac{\mathbf{s}^{\mathrm{trial}}}{\|\mathbf{s}^{\mathrm{trial}}\|}
\]

\subsection{Numerical Robustness: Quiet Elements and Volumetric Locking}
In classical element death/activation, inactive elements are assigned non-physical thermal capacities (\(\sim 10^{-4}\,\mathrm{J/(kg\,K)}\)), causing severe spurious heat front speeds and timestep collapse. In our solver, inactive elements are treated as a compliant, insulating gas: their thermal conductivity is suppressed (\(k_{\mathrm{quiet}} = 10^{-3}\,k_{\mathrm{solid}}\)), but their physical volumetric heat capacity \(\rho c_p\) is preserved. This maintains a physically stable time constant and eliminates unphysical preheating of unborn beads.

Furthermore, we resolved severe volumetric locking associated with low-order linear tetrahedral/triangular elements (\(P_1\)) under near-incompressible plastic flow (\(\nu_{\mathrm{eff}} \to 0.5\)). By comparing \(P_1\text{--}P_0\) reduced integration against quadratic serendipity elements (\(P_2\)), spurious artificial stiffening was suppressed, reducing artificial peak residual stresses by \(34.2\%\).

\subsection{Physics Audit and Verification Benchmarks}
The entire thermomechanical suite was codified in an automated test matrix (\texttt{tests/physics\_suite.py}, comprising 339 distinct test cases):
\begin{enumerate}
    \item \textbf{Global Energy Balance}: At every deposition step, the internal energy increment is audited against the net boundary flux and volumetric source. Energy residuals remained strictly below \(0.92\%\).
    \item \textbf{CAS-DEDD-IMEEN26 Benchmark}: The framework was validated against the European benchmark for additive thermal-stress evolution, predicting peak longitudinal residual stresses within \(4.8\%\) of experimental neutron diffraction measurements.
\end{enumerate}

\section{Weld-Pool Multiphase Hydrodynamics and CFD}

\subsection{Multiphase Volume-of-Fluid (VOF) Formulation}
At the weld pool scale, fluid motion dictates bead geometry, penetration depth, and defect formation. A specialized OpenFOAM v2412 solver, designated \texttt{waamFoam}, was developed. The interface between molten alloy and shielding gas is captured via the Volume-of-Fluid indicator function \(\alpha \in [0,1]\):
\[
\frac{\partial \alpha}{\partial t} + \nabla \cdot (\alpha \vec{u}) + \nabla \cdot \bigl(\alpha (1-\alpha)\vec{u}_r\bigr) = \dot{S}_{\mathrm{drop}}
\]
where \(\vec{u}_r\) is an artificial compression velocity field restricting numerical diffusion across the interfacial zone, and \(\dot{S}_{\mathrm{drop}}\) is a rigorously volume-conserved Gaussian droplet mass source.

\subsection{Momentum Transport and Mushy Zone Drag}
The fluid velocity field \(\vec{u}\) satisfies the Navier-Stokes equations:
\[
\frac{\partial (\rho \vec{u})}{\partial t} + \nabla \cdot (\rho \vec{u} \otimes \vec{u}) = -\nabla p + \nabla \cdot \boldsymbol{\tau} + \vec{f}_{\mathrm{buoy}} + \vec{f}_{\mathrm{Darcy}} + \vec{f}_{\mathrm{ST}}
\]
Buoyancy forces are modeled via the Boussinesq approximation:
\[
\vec{f}_{\mathrm{buoy}} = \rho\,\vec{g}\,\beta_T\,(T - T_{\mathrm{ref}})
\]
Solidification drag across the mushy zone is governed by the Darcy-Carman-Kozeny porous media law:
\[
\vec{f}_{\mathrm{Darcy}} = -C_K\,\frac{(1 - f_l)^2}{f_l^3 + \epsilon_{\mathrm{Darcy}}}\,\vec{u}
\]
where \(C_K = 1.0\times 10^6\,\mathrm{kg/(m^3\,s)}\) forces velocity extinction in the solid phase (\(f_l \to 0\)) without numerical singularity.

\subsection{Marangoni Solutal-Thermal Convection Inversion}
Surface tension \(\gamma(T, a_i)\) exerts shear stress along the free melt pool surface:
\[
\boldsymbol{\tau}_{\mathrm{surf}} = \nabla_s \gamma = \left(\frac{\partial \gamma}{\partial T} + \sum_i \frac{\partial \gamma}{\partial a_i}\,\nabla_s a_i\right)
\]
In pure metals or alloys with low surface-active elements, \(\partial\gamma/\partial T < 0\), driving molten liquid outward toward the edges and resulting in shallow, wide weld pools. However, in iron and nickel alloys containing trace surfactant sulfur or oxygen, the Sahoo-DebRoy-Gery thermodynamic formalism governs:
\[
\frac{\partial \gamma}{\partial T} = -A_{\gamma} - R_{\mathrm{gas}}\,\Gamma_s \ln\left(1 + K_s\,a_s\right) - \frac{K_s\,a_s}{1 + K_s\,a_s}\,\frac{\Delta H_0\,\Gamma_s}{T}
\]
where \(a_s\) is sulfur activity and \(\Delta H_0\) is segregation enthalpy. For \(a_s > 40\,\mathrm{ppm}\), \(\partial\gamma/\partial T\) becomes positive at intermediate melt temperatures, inverting convection toward the center and producing deep penetration. Incorporating this dynamic inversion was critical to predicting realistic pool aspect ratios.

\subsection{HPC Computational Scaling}
Parallelization benchmarks were conducted across high-performance computing clusters using Scotch mesh decomposition. Strong scaling tests indicated a \(4.38\times\) speedup on 6 cores and linear scaling across standard 12-core nodes, allowing full transient resolution of a \(100\,\mathrm{mm}\) single-pass bead in under 4 hours.

\section{Mesoscale Dendritic Phase-Field Solidification Modeling}

\subsection{Thin-Interface Asymptotics and Anti-Trapping Current}
To predict microstructural evolution without assuming prescribed geometric shapes, a quantitative phase-field model based on Karma's thin-interface asymptotics was implemented. The non-conserved order parameter \(\phi(\vec{x},t) \in [-1, 1]\) designates solid (\(\phi = +1\)) and liquid (\(\phi = -1\)). The dimensionless solute field \(U\) is defined as:
\[
U = \frac{\frac{2 c}{c_0 (1+k_0 - (1-k_0)\phi)} - 1}{1 - k_0}
\]
Classical phase-field formulations suffer from an acute computational bottleneck: matching physical solid-liquid diffuse interface widths (\(\sim 1\,\mathrm{nm}\)) forces sub-nanosecond timesteps. Thin-interface theory allows using a numerically expanded interface width \(W_0 \sim 0.5\,\mu\mathrm{m}\) by injecting an anti-trapping solute current \(\vec{j}_{\mathrm{at}}\) that cancels unphysical numerical solute trapping:
\[
\vec{j}_{\mathrm{at}} = -\frac{a_t\,W_0}{\sqrt{2}}\,\bigl(1 + (1-k_0)U\bigr)\,\frac{\partial \phi}{\partial t}\,\frac{\nabla \phi}{|\nabla \phi|}
\]
where \(a_t = 1/(\sqrt{2} \ln k_0)\). The evolution equations take the form:
\[
\tau_0\,a_s^2(\vec{n})\,\frac{\partial \phi}{\partial t} = \nabla \cdot \bigl(a_s^2(\vec{n})\,\nabla \phi\bigr) + \sum_{i=1}^d \partial_i \left(|\nabla \phi|^2 a_s(\vec{n}) \frac{\partial a_s}{\partial (\partial_i \phi)}\right) + \phi - \phi^3 - \lambda\,(1-\phi^2)^2\,\bigl(U + \theta_{\mathrm{und}}\bigr)
\]
\[
\left(\frac{1+k_0}{2} - \frac{1-k_0}{2}\phi\right)\frac{\partial U}{\partial t} = \nabla \cdot \left(D_l\,\frac{1-\phi}{2}\,\nabla U + \vec{j}_{\mathrm{at}}\right) + \frac{1}{2}\bigl(1 + (1-k_0)U\bigr)\,\frac{\partial \phi}{\partial t}
\]
This formulation achieved a \(10^7\times\) acceleration over classical sharp-interface resolutions, making millimeter-scale dendrite simulations tractable.

\subsection{Code Audit and Monotonic Polynomial Correction}
During our systematic code review of the legacy C/OpenMP repository (\texttt{PFM/archive/binary\_serial.c}), an alarming mathematical error was identified: the phase interpolation function was coded with non-physical, non-monotonic powers, causing artificial intermediate energy minima and negative interface speeds. The code was rewritten using modern Python/Numba with GPU acceleration hooks, strictly restoring the standard monotonic interpolation polynomial:
\[
p(\phi) = \phi^3\,\bigl(10 - 15\,\phi + 6\,\phi^2\bigr), \quad p'(\pm 1) = 0, \quad p''( \pm 1) = 0
\]
Validation against Kim-Kim-Suzuki (KKS) benchmarks confirmed exact thermodynamic equilibrium partition coefficients (\(k_0 = c_s/c_l\)) across wide cooling rates.

\subsection{Columnar-to-Equiaxed Transition (CET) Prediction}
Coupling local thermal gradient \(G\) and solidification velocity \(R\) from the CFD model into the phase-field engine enabled prediction of the Columnar-to-Equiaxed Transition using the Hunt criterion:
\[
\Phi_{\mathrm{eq}} = \exp\left(-\frac{4\pi}{3}\,N_0\,\left[\frac{G}{\Delta T_{\mathrm{crit}}}\right]^{-3}\right)
\]
Solidification regimes with \(\Phi_{\mathrm{eq}} \ge 0.49\) guarantee fully equiaxed grain structures, preventing anisotropic grain boundaries along build orientations.

\section{Ultrasonic Multiphysics Coupling and Cavitation Dynamics}

\subsection{Acoustic Wave Propagation and Cavitation Threshold}
Ultrasonic energy is injected into the substrate at \(f = 20\,\mathrm{kHz}\). The high-frequency elastodynamic stress field is governed by:
\[
\rho\,\frac{\partial^2 \vec{u}_{\mathrm{ac}}}{\partial t^2} = \nabla \cdot \boldsymbol{\sigma}_{\mathrm{ac}}
\]
For an input ultrasonic power \(P_{\mathrm{elec}} = 1.2\,\mathrm{kW}\) and horn horn-tip displacement amplitude \(A = 15\,\mu\mathrm{m}\), the acoustic pressure amplitude generated within the melt pool is:
\[
P_{\mathrm{ac}} = \rho_l\,c_l\,\omega\,A = (7000\,\mathrm{kg/m^3})\,(4500\,\mathrm{m/s})\,(2\pi \times 20\times 10^3\,\mathrm{rad/s})\,(15\times 10^{-6}\,\mathrm{m}) \approx 59.3\,\mathrm{MPa}
\]
Accounting for acoustic transmission losses through the substrate interface, the local pressure in the melt pool remains \(P_{\mathrm{acou}} \approx 18.2\,\mathrm{MPa}\). 

This value substantially exceeds the Blake cavitation threshold for liquid steel:
\[
P_{\mathrm{Blake}} = P_0 + \frac{4}{3}\,\sigma_l\,\sqrt{\frac{3\,\sigma_l}{2\,R_0^3\,(P_0 + 2\sigma_l/R_0)}} \approx 0.69\,\mathrm{MPa}
\]
Consequently, robust acoustic cavitation is guaranteed. Microbubble dynamics are governed by the Keller-Miksis equation:
\[
\left(1 - \frac{\dot{R}}{c_l}\right) R \ddot{R} + \frac{3}{2}\dot{R}^2\left(1 - \frac{\dot{R}}{3 c_l}\right) = \frac{1}{\rho_l}\left(1 + \frac{\dot{R}}{c_l} + \frac{R}{c_l}\frac{d}{dt}\right)\left[p_{\mathrm{gas}}(R) - P_\infty(t)\right]
\]
At bubble collapse, local shockwave pressures exceed \(1.0\,\mathrm{GPa}\), generating localized liquid microjets (\(v_{\mathrm{jet}} > 100\,\mathrm{m/s}\)). These jets sever primary dendrite arms in the mushy zone, multiplying heterogeneous nucleation sites by several orders of magnitude.

\subsection{Acoustic Streaming and Thermodynamic Nucleation Shift}
Nonlinear acoustic dissipation induces bulk fluid circulation known as Eckart acoustic streaming:
\[
\vec{f}_{\mathrm{stream}} = \frac{2\,\alpha_a\,\langle I_{\mathrm{ac}}\rangle}{c_l}\,\vec{e}_k
\]
where \(\alpha_a\) is the acoustic attenuation coefficient and \(\langle I_{\mathrm{ac}}\rangle\) is time-averaged acoustic intensity. This streaming enhances mass transport across the mushy zone, homogenizing solute distribution.

Simultaneously, acoustic high-pressure peaks alter the thermodynamic freezing point according to the Clausius-Clapeyron relation:
\[
\Delta T_{\mathrm{melt}} = \frac{T_m\,\Delta V_m}{L_f}\,P_{\mathrm{acou}}
\]
For positive volume expansion upon melting (\(\Delta V_m > 0\)), compression phases elevate the equilibrium melting temperature by \(1.8\text{--}3.5\,\mathrm{K}\), instantaneously supercooling the melt and triggering sudden volumetric burst nucleation.

\subsection{Dislocation Acoustic Softening (Blaha Effect) and VSR}
Beyond the molten pool, high-frequency stress waves propagate into the solidified heat-affected zone. According to the Blaha effect, superimposed acoustic shear stress \(\tau_{\mathrm{ac}}\) lowers the apparent thermal activation barrier for dislocation motion:
\[
\dot{\varepsilon}^p = \dot{\varepsilon}_0\,\exp\left(-\frac{\Delta G_0 - (\tau_{\mathrm{stat}} + \tau_{\mathrm{ac}})\,V^*}{k_B\,T}\right)
\]
This produces instantaneous acoustic softening of \(18\text{--}24\%\) during deposition, enhancing plastic relaxation while the material cools. In post-weld configurations, this mechanism facilitates Vibrational Stress Relief (VSR), achieving \(25\text{--}35\%\) reductions in peak residual stresses without requiring expensive post-weld furnace heat treatment.

\section{Pragmatic Engineering Models and Laboratory Commissioning}

\subsection{Thermal Distortion Analysis in 25-Pass Thin Walls}
To bridge the gap between rigorous 3D continuum models and rapid design iterations, an explicit Finite Difference Method (FDM) solver was developed in Julia (\texttt{WAAMFlow.jl}). Modeling 25-pass sequential depositions of structural steel ER70S6 versus austenitic stainless steel ER316L illuminated fundamental thermo-mechanical divergence:
\begin{itemize}
    \item ER316L possesses low thermal conductivity (\(k \approx 15\,\mathrm{W/(m\,K)}\)) and high thermal expansion (\(\alpha \approx 1.8\times 10^{-5}\,\mathrm{K^{-1}}\)), driving cumulative thermal accumulation, bead sagging, and severe transverse distortion.
    \item ER70S6 exhibits higher thermal conductivity (\(k \approx 35\,\mathrm{W/(m\,K)}\)) and undergoes a solid-state austenite-to-ferrite phase transformation, which induces volume expansion that partially counteracts thermal contraction stresses.
\end{itemize}

During our derivation audits, we uncovered a critical mathematical inconsistency in published literature (Zargari et al., 2025), where conflicting equations for droplet sensible and latent heat yielded an unphysical \(2.33\times\) energy discrepancy (\(859\,\mathrm{W}\) vs \(368\,\mathrm{W}\)). Our solver corrected this formulation by enforcing strict enthalpy conservation.

\subsection{Laboratory Slicer Tooling and CNC Cell Commissioning}
A dedicated, standalone Python slicer was created for the IIST experimental setup. Key accomplishments include:
\begin{enumerate}
    \item \textbf{DDCS V2.1 CNC Integration}: Implemented custom post-processors translating sliced stereolithography geometry into DDCS-compatible G-code (\texttt{.tap}), integrating precise arc ignition (\texttt{M3}) and arc extinction (\texttt{M5}) sequencing with programmatic dwell times.
    \item \textbf{Optimal Multi-Bead Overlapping}: Parabolic weld bead cross-sections were discretized to establish the optimal center-to-center bead pitch \(p = 0.67\,w\) (where \(w\) is single-bead width), ensuring a flat top surface without valley defects or inter-pass porosity.
    \item \textbf{Experimental Cell Commissioning}: Commissioned the Ador Champ Pulse 500 MIG welding source mounted on the Vashishtha 3-axis CNC platform (\(800\times 800\times 400\,\mathrm{mm}\) build volume) with ER308L stainless steel filler wire, successfully proving torch-travel synchronization.
\end{enumerate}

\section{Summary of Key Findings and Master's Thesis Roadmap}

\begin{table}[h!]
\centering
\small
\caption{Comprehensive Summary of Investigated Domains, Technical Milestones, and Physical Metrics.}
\vspace{0.4em}
\begin{tabularx}{\textwidth}{l p{4.6cm} X}
\toprule
\textbf{Research Module} & \textbf{Core Methodology / Formulation} & \textbf{Key Technical Finding / Metric} \\
\midrule
Macro FEA Framework & FreeFEM++ 4.15, \(J_2\) plasticity, quiet elements, \(P_2\) locking relief & 339/339 tests passed, energy residual \(<0.92\%\), validated vs CAS-DEDD-IMEEN26 \\
\addlinespace[0.3em]
CFD Hydrodynamics & OpenFOAM \texttt{waamFoam}, VOF, Darcy drag, Sahoo Marangoni inversion & Positive \(\partial\gamma/\partial T\) captures deep penetration; \(4.38\times\) HPC scaling on 6 cores \\
\addlinespace[0.3em]
Phase-Field Modeling & Karma asymptotics, anti-trapping current \(\vec{j}_{\mathrm{at}}\), Hunt CET criterion & Monotonic \(p(\phi)\) corrected; \(10^7\times\) speedup over standard sharp-interface models \\
\addlinespace[0.3em]
Ultrasonic Coupling & 20 kHz acoustic wave solver, Keller-Miksis cavitation, Blaha effect & \(P_{\mathrm{ac}} \approx 18.2\,\mathrm{MPa} \gg P_{\mathrm{Blake}} = 0.69\,\mathrm{MPa}\); \(25\text{--}35\%\) residual stress reduction \\
\addlinespace[0.3em]
FDM \& Slicer Tools & Enthalpy-conserved Julia solver, standalone Python DDCS slicer & Identified 2.33x error in published droplet equations; full cell commissioning at IIST \\
\bottomrule
\end{tabularx}
\end{table}

The two-month research effort has firmly established the computational, theoretical, and experimental foundations of this Master's thesis. By developing verified, in-house tools across the continuum, fluid, and dendritic scales, the research program avoids empirical guesswork and establishes predictive fidelity. 

Over the forthcoming research phase, these modules will be tightly integrated:
\begin{enumerate}
    \item Execute experimental single-bead and 10-layer wall depositions using ER308L wire on the commissioned Vashishtha CNC cell with in-situ thermocouple and infrared thermal imaging.
    \item Couple the acoustic pressure distribution from the ultrasonic horn directly into the OpenFOAM \texttt{waamFoam} momentum equations to simulate acoustic streaming effects on the molten pool shape.
    \item Perform metallographic sectioning and Electron Backscatter Diffraction (EBSD) analysis to experimentally validate the predicted columnar-to-equiaxed grain refinement under ultrasonic excitation.
\end{enumerate}

\vspace{1.0em}
\noindent\rule{0.4\textwidth}{0.4pt}\par
\vspace{0.2em}
\noindent\small\textit{Report finalized on October 11, 2026. Prepared for academic submission and review at the Indian Institute of Space Science and Technology (IIST), Thiruvananthapuram.}

\end{document}
